% ===========================================================================
% AIG230 -- Natural Language Processing
% Final Project, Option 1: Transformer-Based Neural Machine Translation
%
% Group 7: Jose Luis Sanchez Noriega, Bikash Subedi
% Instructor: Ellie Azizi -- Seneca Polytechnic
%
% Build:  latexmk -pdf main.tex
%     or: make report      (from the repository root)
%
% Numbers and tables come from reports/build_report_data.py, which reads the
% JSON artefacts under artifacts/results/. Regenerate them before compiling:
%     python reports/build_report_data.py
% ===========================================================================
\documentclass[11pt,a4paper]{article}

% ---------------------------------------------------------------------------
% Shared preamble for the AIG230 final report.
% ---------------------------------------------------------------------------

\usepackage[utf8]{inputenc}
\usepackage[T1]{fontenc}
\usepackage[english]{babel}

% --- Typography ------------------------------------------------------------
% Libertinus is a Linux Libertine derivative: an old-style serif with a large
% x-height that stays readable at 11pt, paired with a matching sans for
% headings. Chosen to echo the typographic feel of distill.pub articles, whose
% figure style this report also follows.
\usepackage{libertinus}
\usepackage{microtype}

\usepackage[a4paper,
            top=2.5cm, bottom=2.5cm, left=2.7cm, right=2.7cm,
            headheight=14pt]{geometry}

% --- Mathematics -----------------------------------------------------------
\usepackage{amsmath}
\usepackage{amssymb}
\usepackage{bm}

% --- Tables and figures ----------------------------------------------------
\usepackage{booktabs}
\usepackage{longtable}
\usepackage{array}
\usepackage{multirow}
\usepackage{graphicx}
\usepackage{float}
\usepackage[font=small,labelfont=bf,justification=raggedright,
            singlelinecheck=false]{caption}
\usepackage{subcaption}

% --- Colour ----------------------------------------------------------------
\usepackage{xcolor}
\definecolor{accentblue}{HTML}{2A78D6}
\definecolor{accentorange}{HTML}{EB6834}
\definecolor{accentgreen}{HTML}{1BAF7A}
\definecolor{inkprimary}{HTML}{0B0B0B}
\definecolor{inksecondary}{HTML}{52514E}
\definecolor{inkmuted}{HTML}{898781}
\definecolor{rule}{HTML}{E1E0D9}
\definecolor{surface}{HTML}{FCFCFB}

% --- Boxes -----------------------------------------------------------------
\usepackage[most]{tcolorbox}

\newtcolorbox{keypoint}[1][]{
  enhanced, breakable,
  colback=surface, colframe=accentblue,
  boxrule=0pt, leftrule=2.5pt,
  arc=1pt, left=8pt, right=8pt, top=6pt, bottom=6pt,
  fontupper=\small, #1
}

\newtcolorbox{decisionbox}[1][]{
  enhanced, breakable,
  colback=surface, colframe=accentgreen,
  boxrule=0pt, leftrule=2.5pt,
  arc=1pt, left=8pt, right=8pt, top=6pt, bottom=6pt,
  fontupper=\small, #1
}

% Marks a paragraph whose wording depends on final trained numbers.
\newtcolorbox{pending}[1][]{
  enhanced, breakable,
  colback=yellow!6, colframe=accentorange,
  boxrule=0pt, leftrule=2.5pt,
  arc=1pt, left=8pt, right=8pt, top=6pt, bottom=6pt,
  fontupper=\small\itshape, #1
}

% --- Headings --------------------------------------------------------------
\usepackage{titlesec}
\titleformat{\section}{\Large\bfseries\color{inkprimary}}{\thesection}{0.7em}{}
\titleformat{\subsection}{\large\bfseries\color{inkprimary}}{\thesubsection}{0.6em}{}
\titleformat{\subsubsection}{\normalsize\bfseries\color{inksecondary}}
            {\thesubsubsection}{0.6em}{}
\titlespacing*{\section}{0pt}{2.0ex plus 1ex minus .2ex}{1.1ex plus .2ex}

% --- Lists -----------------------------------------------------------------
\usepackage{enumitem}
\setlist{itemsep=2pt, topsep=4pt, parsep=0pt}

% --- Headers ---------------------------------------------------------------
\usepackage{fancyhdr}
\pagestyle{fancy}
\fancyhf{}
\fancyhead[L]{\small\color{inkmuted}AIG230 Final Project --- Group 7}
\fancyhead[R]{\small\color{inkmuted}\thepage}
\renewcommand{\headrulewidth}{0.4pt}
\renewcommand{\headrule}{\hbox to\headwidth{\color{rule}\leaders\hrule height \headrulewidth\hfill}}

% --- Code ------------------------------------------------------------------
\usepackage{listings}
\lstset{
  basicstyle=\ttfamily\footnotesize,
  keywordstyle=\color{accentblue}\bfseries,
  commentstyle=\color{inkmuted}\itshape,
  stringstyle=\color{accentgreen},
  showstringspaces=false,
  breaklines=true,
  frame=leftline,
  framesep=6pt,
  rulecolor=\color{rule},
  xleftmargin=10pt,
  numbers=none,
  language=Python,
}

% --- Links -----------------------------------------------------------------
\usepackage[hidelinks,
            colorlinks=true,
            linkcolor=accentblue,
            citecolor=accentblue,
            urlcolor=accentblue]{hyperref}
\usepackage{url}

% --- Spacing ---------------------------------------------------------------
\setlength{\parskip}{0.55em}
\setlength{\parindent}{0pt}
\renewcommand{\arraystretch}{1.15}

% --- Convenience -----------------------------------------------------------
\newcommand{\code}[1]{\texttt{\small #1}}
\newcommand{\dmodel}{d_{\mathrm{model}}}
\newcommand{\dff}{d_{\mathrm{ff}}}
\newcommand{\dk}{d_k}
\newcommand{\entos}{EN\,$\rightarrow$\,ES}
\newcommand{\estoen}{ES\,$\rightarrow$\,EN}

% ---------------------------------------------------------------------------
% Fallbacks so the document compiles before any model has been trained.
% `generated/macros.tex` overrides whichever of these it defines.
% ---------------------------------------------------------------------------
\newcommand{\definefallbacks}{%
  \providecommand{\HasResults}{false}%
  \providecommand{\PrimaryRun}{BPE 16k, learned embeddings}%
  % corpus
  \providecommand{\TatoebaEnglishSentences}{--}%
  \providecommand{\TatoebaSpanishSentences}{--}%
  \providecommand{\TatoebaLinksScanned}{--}%
  \providecommand{\PairsRaw}{--}%
  \providecommand{\PairsClean}{--}%
  \providecommand{\RetentionPercent}{--}%
  \providecommand{\PairsTrain}{--}%
  \providecommand{\PairsValidation}{--}%
  \providecommand{\PairsTest}{--}%
  \providecommand{\Components}{--}%
  \providecommand{\Leakage}{--}%
  \providecommand{\LargestComponent}{--}%
  \providecommand{\TypesEn}{--}%
  \providecommand{\TypesEs}{--}%
  \providecommand{\TypeRatio}{--}%
  \providecommand{\TokensEn}{--}%
  \providecommand{\TokensEs}{--}%
  \providecommand{\HapaxEnPercent}{--}%
  \providecommand{\HapaxEsPercent}{--}%
  \providecommand{\OovEnPercent}{--}%
  \providecommand{\OovEsPercent}{--}%
  \providecommand{\FertilityEn}{--}%
  \providecommand{\FertilityEs}{--}%
  \providecommand{\SubwordVocab}{--}%
  \providecommand{\WordVocab}{--}%
  \providecommand{\MuseCoveragePercent}{--}%
  \providecommand{\MuseCovered}{--}%
  \providecommand{\MuseRandom}{--}%
  % per-run parameters, epochs and scores
  \providecommand{\ParamsBpeScratch}{--}%
  \providecommand{\ParamsWordRandom}{--}%
  \providecommand{\ParamsWordMuse}{--}%
  \providecommand{\ParamsLstmBaseline}{--}%
  \providecommand{\EpochsBpeScratch}{--}%
  \providecommand{\EpochsWordRandom}{--}%
  \providecommand{\EpochsWordMuse}{--}%
  \providecommand{\EpochsLstmBaseline}{--}%
  \providecommand{\BleuBpeScratchEnEs}{--}%
  \providecommand{\BleuBpeScratchEsEn}{--}%
  \providecommand{\BleuBpeScratchMean}{--}%
  \providecommand{\BleuWordRandomEnEs}{--}%
  \providecommand{\BleuWordRandomEsEn}{--}%
  \providecommand{\BleuWordRandomMean}{--}%
  \providecommand{\BleuWordMuseEnEs}{--}%
  \providecommand{\BleuWordMuseEsEn}{--}%
  \providecommand{\BleuWordMuseMean}{--}%
  \providecommand{\BleuLstmBaselineEnEs}{--}%
  \providecommand{\BleuLstmBaselineEsEn}{--}%
  \providecommand{\BleuLstmBaselineMean}{--}%
  \providecommand{\ChrfBpeScratchEnEs}{--}%
  \providecommand{\ChrfBpeScratchEsEn}{--}%
}

% Include a figure only if the file exists, so a partially-populated
% figures directory never breaks the build.
%
% The figure names produced by nmt.viz contain underscores (`data_construction`),
% which TeX reads as subscript markers and which would otherwise raise
% "Missing $ inserted" the moment the name reaches \includegraphics.
% \detokenize turns the argument into plain characters before that happens.
% \IfFileExists searches the current directory and TEXINPUTS -- it does *not*
% consult \graphicspath. Passing a bare name therefore always took the "not
% found" branch and silently rendered a placeholder box for every figure, even
% though the PDFs were present. The directory is given explicitly instead;
% each document sets \figuredir to point at reports/figures/.
\providecommand{\figuredir}{../figures/}

\newcommand{\maybefigure}[4]{%
  \IfFileExists{\figuredir\detokenize{#1}.pdf}{%
    \begin{figure}[htbp]
      \centering
      \includegraphics[width=#2\linewidth]{\figuredir\detokenize{#1}}
      \caption{#3}
      \label{#4}
    \end{figure}%
  }{%
    \begin{figure}[htbp]
      \centering
      \fbox{\parbox{0.8\linewidth}{\centering\small\itshape
        Figure not generated yet: \code{\detokenize{#1}.pdf}\\[2pt]
        Run \code{python -m nmt.viz.make\_figures}.}}
      \caption{#3}
      \label{#4}
    \end{figure}%
  }%
}

% Order matters: the generated file uses \newcommand and must be read first,
% then \definefallbacks fills in (with \providecommand) only those macros the
% generator had no data for. Reversing this would make every real value clash
% with its own placeholder.
\IfFileExists{generated/macros.tex}{% Generated by reports/build_report_data.py -- do not edit by hand.
% Every number quoted inline in the report comes from here, so the
% prose cannot drift out of step with the artefacts in the repo.

\newcommand{\BleuBpeScratchEnEs}{49.99}
\newcommand{\BleuBpeScratchEsEn}{56.83}
\newcommand{\BleuBpeScratchMean}{53.41}
\newcommand{\BleuLstmBaselineEnEs}{46.15}
\newcommand{\BleuLstmBaselineEsEn}{49.73}
\newcommand{\BleuLstmBaselineMean}{47.94}
\newcommand{\BleuWordMuseEnEs}{18.15}
\newcommand{\BleuWordMuseEsEn}{23.57}
\newcommand{\BleuWordMuseMean}{20.86}
\newcommand{\BleuWordRandomEnEs}{45.42}
\newcommand{\BleuWordRandomEsEn}{52.62}
\newcommand{\BleuWordRandomMean}{49.02}
\newcommand{\ChrfBpeScratchEnEs}{68.50}
\newcommand{\ChrfBpeScratchEsEn}{71.43}
\newcommand{\ChrfLstmBaselineEnEs}{65.71}
\newcommand{\ChrfLstmBaselineEsEn}{66.37}
\newcommand{\ChrfWordMuseEnEs}{42.35}
\newcommand{\ChrfWordMuseEsEn}{43.62}
\newcommand{\ChrfWordRandomEnEs}{65.57}
\newcommand{\ChrfWordRandomEsEn}{68.22}
\newcommand{\Components}{222,950}
\newcommand{\EpochsBpeScratch}{20}
\newcommand{\EpochsLstmBaseline}{9}
\newcommand{\EpochsWordMuse}{6}
\newcommand{\EpochsWordRandom}{19}
\newcommand{\ErrorBpeScratchCopiedSourceEnEs}{0.2}
\newcommand{\ErrorBpeScratchCopiedSourceEsEn}{0.1}
\newcommand{\ErrorBpeScratchNoContentOverlapEnEs}{3.8}
\newcommand{\ErrorBpeScratchNoContentOverlapEsEn}{2.3}
\newcommand{\ErrorBpeScratchNumberMismatchEnEs}{0.2}
\newcommand{\ErrorBpeScratchNumberMismatchEsEn}{0.4}
\newcommand{\ErrorBpeScratchOtherEnEs}{94.7}
\newcommand{\ErrorBpeScratchOtherEsEn}{96.3}
\newcommand{\ErrorBpeScratchOverGenerationEnEs}{0.4}
\newcommand{\ErrorBpeScratchOverGenerationEsEn}{0.3}
\newcommand{\ErrorBpeScratchRepetitionEnEs}{0.6}
\newcommand{\ErrorBpeScratchRepetitionEsEn}{0.6}
\newcommand{\ErrorBpeScratchTruncationEnEs}{0.4}
\newcommand{\ErrorBpeScratchTruncationEsEn}{0.4}
\newcommand{\FertilityEn}{1.35}
\newcommand{\FertilityEs}{1.35}
\newcommand{\HapaxEnPercent}{39.6}
\newcommand{\HapaxEsPercent}{42.5}
\newcommand{\HasResults}{true}
\newcommand{\LargestComponent}{169}
\newcommand{\Leakage}{0}
\newcommand{\MeanSentenceBleuBpeScratchEnEs}{49.82}
\newcommand{\MeanSentenceBleuBpeScratchEsEn}{57.02}
\newcommand{\MuseCoveragePercent}{92.2}
\newcommand{\MuseCovered}{29,511}
\newcommand{\MuseRandom}{2,483}
\newcommand{\OovEnPercent}{1.59}
\newcommand{\OovEsPercent}{2.75}
\newcommand{\PairsClean}{282,785}
\newcommand{\PairsRaw}{282,902}
\newcommand{\PairsTest}{5,656}
\newcommand{\PairsTrain}{271,473}
\newcommand{\PairsValidation}{5,656}
\newcommand{\ParamsBpeScratch}{37,635,712}
\newcommand{\ParamsLstmBaseline}{38,921,920}
\newcommand{\ParamsWordMuse}{39,366,912}
\newcommand{\ParamsWordRandom}{39,366,912}
\newcommand{\PrimaryRun}{BPE 16k, learned embeddings}
\newcommand{\RetentionPercent}{99.96}
\newcommand{\SubwordVocab}{16,000}
\newcommand{\TatoebaEnglishSentences}{2,033,133}
\newcommand{\TatoebaLinksScanned}{28,361,472}
\newcommand{\TatoebaSpanishSentences}{441,233}
\newcommand{\TokensEn}{2,161,611}
\newcommand{\TokensEs}{2,145,680}
\newcommand{\TypeRatio}{1.61}
\newcommand{\TypesEn}{35,066}
\newcommand{\TypesEs}{56,519}
\newcommand{\WordVocab}{32,000}
}{}
\definefallbacks

\graphicspath{{../figures/}{figures/}}

\title{\vspace{-1.5cm}\bfseries Transformer-Based Neural Machine Translation\\[0.25em]
  \large\mdseries A single bidirectional English\,$\leftrightarrow$\,Spanish model}
\author{}
\date{}

\begin{document}

% ---------------------------------------------------------------------------
\begin{titlepage}
\centering
\vspace*{2.2cm}

{\color{inkmuted}\large Seneca Polytechnic \\[0.3em]
AIG230 --- Natural Language Processing \\[0.3em]
Postgraduate Certificate in Artificial Intelligence}

\vspace{1.8cm}
{\color{rule}\rule{0.72\linewidth}{1.2pt}}
\vspace{1.2cm}

{\Huge\bfseries Transformer-Based\\[0.35em] Neural Machine Translation\par}

\vspace{0.9cm}
{\LARGE A single bidirectional English\,$\leftrightarrow$\,Spanish model\par}

\vspace{0.6cm}
{\color{rule}\rule{0.72\linewidth}{1.2pt}}

\vspace{2.2cm}
{\large\bfseries Group 7\par}
\vspace{0.5em}
{\large Jose Luis Sanchez Noriega \\[0.35em] Bikash Subedi\par}

\vspace{2.0cm}
{\large Instructor: Ellie Azizi\par}

\vfill
{\color{inksecondary}\small
Final Project --- Option 1 \\[0.4em]
Corpus: the Tatoeba Project (English--Spanish) \\[0.4em]
Framework: PyTorch \\[0.8em]
\today}
\end{titlepage}

% ---------------------------------------------------------------------------
\section*{Abstract}
\addcontentsline{toc}{section}{Abstract}

We build a complete neural machine translation system for English and Spanish
using a transformer encoder--decoder implemented from primitives in PyTorch.
The distinguishing design choice is that \emph{one} set of weights serves
\emph{both} translation directions: source and target draw on a single joint
subword vocabulary, the embedding matrix is shared three ways between the
encoder, the decoder and the output projection, and a reserved direction tag
prepended to the source selects the output language. This halves the parameter
budget relative to two separate models while doubling the supervision each
parameter receives.

The corpus is reconstructed from the official Tatoeba exports by joining
\TatoebaEnglishSentences{} English sentences and \TatoebaSpanishSentences{}
Spanish sentences through \TatoebaLinksScanned{} translation links, yielding
\PairsRaw{} aligned pairs. We show that the standard practice of splitting such
a corpus uniformly at random leaks between train and test --- Tatoeba is a
translation \emph{graph}, not a list of independent pairs --- and partition
connected components instead, which removes the leak entirely.

We report BLEU and chrF2 in both directions against two baselines: a
capacity-matched recurrent sequence-to-sequence model with Bahdanau attention,
and an ablation isolating the effect of initialising embeddings from MUSE
cross-lingual vectors. An automated error analysis classifies every test
translation into named failure modes and surfaces the twenty worst cases for
manual inspection. The trained model is served through a Gradio application
that translates either direction interactively and displays the
cross-attention alignment behind each output.

\vspace{0.6em}
\begin{keypoint}
\textbf{Reproducibility.} Every number, table and figure in this report is
generated from JSON artefacts written by the pipeline
(\code{reports/build\_report\_data.py} and \code{nmt.viz.make\_figures}).
Nothing is transcribed by hand, so the document cannot drift out of step with
the code and checkpoints in the repository.
\end{keypoint}

\clearpage
\tableofcontents
\clearpage

% ---------------------------------------------------------------------------
\section{Introduction}

\subsection{The task}

Machine translation maps a sentence in a source language to a sentence in a
target language that preserves its meaning. It is a \emph{sequence-to-sequence}
problem with three properties that make it hard and that shape every decision
in this project:

\begin{itemize}
  \item \textbf{The output length is not determined by the input length.}
    ``I'm tired'' is two words in English and two in Spanish (``Estoy
    cansado''), but ``I'll see you tomorrow'' is five and ``Te veré mañana'' is
    three. The model must decide when to stop.
  \item \textbf{The alignment is not monotonic.} English places adjectives
    before nouns and Spanish generally after: ``the red house'' becomes ``la
    casa roja''. A model that could only attend to source positions in order
    would be unable to produce this.
  \item \textbf{The target is not unique.} ``I'm tired'' is equally well
    ``Estoy cansado'', ``Estoy cansada'' or ``Tengo sueño''. Any loss that
    demands probability~1 on one particular reference is asking for something
    false --- a fact that directly motivates the label smoothing of
    Section~\ref{sec:training}.
\end{itemize}

\subsection{Why one model for both directions}

The obvious design is two models, one per direction. We chose a single
bidirectional model instead, following the multilingual tagging approach of
Johnson et al.~\cite{johnson2017}. Two reserved tokens, \code{<2en>} and
\code{<2es>}, are prepended to the source sentence:

\begin{center}
\begin{tabular}{@{}ll@{}}
\code{<2es> I am tired .} & $\rightarrow$ \quad \code{Estoy cansado .} \\
\code{<2en> Estoy cansado .} & $\rightarrow$ \quad \code{I am tired .} \\
\end{tabular}
\end{center}

Every cleaned sentence pair therefore produces \emph{two} training examples,
one per direction. Three consequences follow, and they are the reason for the
choice:

\begin{enumerate}
  \item \textbf{Twice the supervision per parameter.} The same encoder must
    represent both English and Spanish sentences, and the same decoder must
    generate both. On a corpus of this size --- around a quarter of a million
    pairs, two orders of magnitude smaller than the WMT datasets transformers
    were designed for --- data, not compute, is the binding constraint, and
    doubling the effective corpus is worth more than the capacity given up.
  \item \textbf{A shared vocabulary becomes possible, and then necessary.}
    Since the decoder emits into the same space the encoder reads from, source
    and target must share one inventory. This in turn permits three-way weight
    tying (Section~\ref{sec:architecture}), removing roughly a third of the
    parameters.
  \item \textbf{Cognates share statistical strength.} ``nation''/``nación'',
    ``important''/``importante'' and ``possible''/``posible'' share subword
    units under a joint BPE model, so evidence about one improves the
    representation of the other.
\end{enumerate}

The cost is capacity: one model of a given size must hold both directions.
Section~\ref{sec:evaluation} reports what that costs in BLEU.

\subsection{Contributions}

\begin{enumerate}
  \item A transformer encoder--decoder written from primitives --- attention,
    positional encoding, masking, residual blocks --- rather than assembled
    from \code{torch.nn.Transformer}, with the training loop written by hand as
    the brief requires.
  \item A corpus construction pipeline that rebuilds the bitext from the
    official Tatoeba exports and demonstrates, with a measurement, that
    component-aware splitting is necessary to avoid leakage.
  \item A controlled three-way comparison isolating the effect of pre-trained
    cross-lingual embeddings from the effect of the tokenisation change that
    normally confounds it.
  \item A recurrent baseline matched on data, vocabulary, optimiser and
    parameter count, so the architectural comparison means something.
  \item An automated error analysis that names and counts failure modes across
    the whole test set rather than reporting a single aggregate score.
\end{enumerate}

% ---------------------------------------------------------------------------
\section{Data collection and exploration}
\label{sec:data}

\subsection{Building the corpus from the official exports}

The Tatoeba Project publishes per-language sentence dumps and a single global
table of translation links; it does not publish a ready-made English--Spanish
bitext. We reconstruct one, which keeps the corpus pinned to a dated snapshot
we control and preserves the Tatoeba sentence identifiers so that any example
quoted in this report can be traced back to the public database.

Three files are downloaded (about 172\,MB compressed) and joined in three
passes so that peak memory stays proportional to the two languages of interest
rather than to the size of the link table:

\begin{enumerate}
  \item load \code{eng\_sentences.tsv} into \code{\{id: text\}} ---
    \TatoebaEnglishSentences{} rows;
  \item load \code{spa\_sentences.tsv} --- \TatoebaSpanishSentences{} rows;
  \item stream \code{links.csv} (\TatoebaLinksScanned{} rows), emitting a pair
    whenever the left identifier is English \emph{and} the right one is
    Spanish.
\end{enumerate}

Because Tatoeba stores every link in both orientations, that asymmetric test
yields each English--Spanish pair exactly once. The join produces
\PairsRaw{} aligned pairs.

\maybefigure{data_construction}{1.0}{From the raw Tatoeba exports to the final
splits. Cleaning retains \RetentionPercent\% of pairs: Tatoeba is a curated,
human-checked corpus, so the filters largely confirm its quality rather than
repair it. The splits are made by connected component, and the leakage check
is run on every build.}{fig:construction}

\subsection{What the corpus looks like}

\begin{table}[htbp]
\centering
\caption{Corpus statistics after cleaning. Subword counts use the joint
\SubwordVocab-entry BPE model fitted on the training split only.}
\label{tab:dataset}
\small
\begin{tabular}{lrrrrrrr}
\toprule
& & \multicolumn{2}{c}{unique sentences} & \multicolumn{2}{c}{mean words} & \multicolumn{2}{c}{mean subwords} \\
\cmidrule(lr){3-4}\cmidrule(lr){5-6}\cmidrule(lr){7-8}
Split & Pairs & EN & ES & EN & ES & EN & ES \\
\midrule
Train & 271,473 & 232,777 & 250,057 & 6.8 & 6.6 & 9.1 & 8.8 \\
Validation & 5,656 & 5,026 & 5,330 & 6.9 & 6.7 & 9.4 & 9.1 \\
Test & 5,656 & 5,032 & 5,326 & 6.9 & 6.7 & 9.3 & 9.0 \\
\bottomrule
\end{tabular}

\end{table}

Three characteristics shape the modelling decisions that follow.

\paragraph{The sentences are short.} Both sides average under seven whitespace
tokens, and the 95th percentile is around thirteen. Tatoeba is a corpus of
everyday example sentences contributed by language learners, not of newswire or
parliamentary proceedings. This is convenient --- attention is quadratic in
sequence length, so short sentences make the model cheap to train --- but it
also bounds what the system can be expected to do: it will be weakest on
exactly the long, syntactically complex sentences that are rarest in training.

\paragraph{Spanish has far more surface forms than English.} The training split
contains \TypesEn{} distinct English word types against \TypesEs{} Spanish
ones, from a near-identical number of running tokens (\TokensEn{} and
\TokensEs{}). That is a factor of \TypeRatio{}. The cause is morphology:
Spanish marks person, number, tense and mood on the verb (\emph{hablo, hablas,
habla, hablamos, habláis, hablan}, and that is one tense of one verb), agrees
adjectives and articles for gender and number, and attaches clitic pronouns
directly to infinitives and gerunds (\emph{dármelo}). English does almost none
of this.

\paragraph{The tail is heavy.} \HapaxEnPercent\% of English types and
\HapaxEsPercent\% of Spanish types occur exactly once in training. A word seen
once cannot train a useful embedding.

\maybefigure{data_vocabulary}{1.0}{The case for subword tokenisation. Spanish
contributes \TypeRatio$\times$ as many word types as English (left); a
\WordVocab-entry word vocabulary still leaves \OovEsPercent\% of Spanish test
tokens unknown against \OovEnPercent\% of English ones (right). Subword
tokenisation removes the unknown-word problem entirely, at the cost of longer
sequences.}{fig:vocabulary}

\maybefigure{data_lengths}{1.0}{Sentence-length distributions and the
relationship between the two sides. The strong correlation is what makes a
length-ratio filter a sound way to detect misaligned links.}{fig:lengths}

\subsection{Characteristics that will influence the task}

\begin{decisionbox}
\textbf{Implications carried into the rest of the project.}
\begin{itemize}
  \item \emph{Short sentences} $\Rightarrow$ a maximum length of 64 tokens
    discards a negligible tail; batches can be capped by token count to keep
    memory flat.
  \item \emph{Spanish morphology} $\Rightarrow$ a subword vocabulary, and
    chrF2 reported alongside BLEU because a translation can get a stem right
    and its inflection wrong, which BLEU scores as a total miss.
  \item \emph{Heavy tail} $\Rightarrow$ a modest \SubwordVocab-entry
    vocabulary, so that even the rare half of it is updated often enough to
    learn.
  \item \emph{Many-to-many links} $\Rightarrow$ component-aware splitting
    (Section~\ref{sec:splitting}).
\end{itemize}
\end{decisionbox}

% ---------------------------------------------------------------------------
\section{Preprocessing}
\label{sec:preprocessing}

Every rule below is a modelling decision, so each is stated with its
justification. The guiding principle is \emph{remove noise, preserve meaning}:
normalise away distinctions the model should not spend capacity on, and leave
intact everything a reader would consider part of the sentence.

\subsection{Normalisation}

Applied in this order, because the order matters:

\begin{enumerate}
  \item \textbf{Unicode NFC composition.} Spanish accents arrive both
    pre-composed (\code{ñ} = U+00F1) and decomposed (\code{n} + U+0303). They
    render identically but are different strings; without NFC the tokeniser
    treats \emph{español} as two distinct types.
  \item \textbf{Control-character removal}, before whitespace collapsing, so
    that a stripped control character cannot weld two words together.
  \item \textbf{Typographic folding.} Five kinds of apostrophe, six of quote
    and seven of dash are mapped to one ASCII form each. Tatoeba is
    crowd-sourced, so the same sentence appears with curly or straight quotes
    depending on the contributor's keyboard.
  \item \textbf{Whitespace collapsing}, last, since the previous steps can
    introduce runs of spaces.
\end{enumerate}

\begin{decisionbox}
\textbf{Casing and punctuation are preserved.} Lowercasing shrinks the
vocabulary and typically buys a point of BLEU on a small corpus, but it makes
the system unable to produce correctly-cased output, which would then need a
separate truecasing model at inference. Since the deliverable is an interactive
application whose output a human reads, correct casing is part of the product.
Likewise, Spanish inverted marks (\code{¿}, \code{¡}) are grammatical signal,
and translating punctuation is part of translating.
\end{decisionbox}

\subsection{Filtering}

Pairs are discarded when either side is empty after normalisation, contains no
alphabetic content (``1997.''), contains a URL or e-mail address, contains
characters from an unexpected script (mislabelled rows), falls outside
1--64 tokens or 2--400 characters, is identical on both sides, or has a
length ratio above 3.0 once both sides reach four tokens. The ratio test is
suppressed on very short sentences because ``Yes.''\ against ``Sí, lo haré.''\
is a legitimate 1:3 pair.

Exact duplicate pairs are removed. Near-duplicates are deliberately \emph{not}
removed: Tatoeba's alternative translations of the same sentence are genuine
paraphrases and discarding them would throw away real supervision. What must
not happen is those paraphrases straddling the train/test boundary --- and that
is handled by the split, not by the filter.

Cleaning retains \PairsClean{} of \PairsRaw{} pairs (\RetentionPercent\%). The
high retention is itself a finding worth stating: Tatoeba is human-curated, so
the aggressive cleaning pipelines that web-crawled corpora require are not
needed here. The filters earn their place by \emph{confirming} data quality and
by catching the small number of genuinely broken rows.

\subsection{Splitting without leakage}
\label{sec:splitting}

\begin{keypoint}
Tatoeba is a translation \emph{graph}, not a list of independent pairs. One
English sentence typically links to several Spanish sentences, and each of
those may link to further English sentences. Splitting the pair list uniformly
at random therefore leaks: ``I'm tired.'' / ``Estoy cansado.'' lands in
training while ``I'm tired.'' / ``Estoy cansada.'' lands in test, and the
reported BLEU measures memorisation rather than translation.
\end{keypoint}

The fix is to split \emph{connected components}. Treat every distinct sentence
as a node and every link as an edge, merge them with a union--find structure,
and assign whole components to a split. Sentences are namespaced by language so
that an English string equal to a Spanish string (proper nouns, ``No.'') does
not spuriously merge two unrelated components.

This yields \Components{} components over \PairsClean{} pairs, the largest
containing \LargestComponent{} pairs. Components contributing more than four
pairs are forced into training: a forty-pair component landing in test would
fill the evaluation set with paraphrases of one sentence and make BLEU hostage
to whether the model happened to know it.

A leakage assertion runs on every build and is recorded in the split report; it
counts sentence strings shared between splits on either side. The current value
is \textbf{\Leakage}.

\subsection{Tokenisation}

Two tokenisers are fitted --- on the \emph{training split only}, since fitting
on the full corpus would leak test-set surface forms into the vocabulary and
flatter the reported out-of-vocabulary rate.

\paragraph{Joint subword (primary).} A SentencePiece BPE model with
\SubwordVocab{} entries, fitted on the concatenation of English and Spanish
training text. Sharing one vocabulary is what makes the bidirectional model
possible. Being open-vocabulary, it never emits \code{<unk>}. Fertility is
\FertilityEn{} subword tokens per English word and \FertilityEs{} per Spanish
word --- the price paid for eliminating unknown words.

\paragraph{Word level (for the embedding experiment).} A frequency-truncated
vocabulary of \WordVocab{} entries with singletons dropped. It exists because
pre-trained word vectors are distributed as word-level tables: to initialise
embeddings from MUSE, the units the model embeds must \emph{be} words. The cost
is a closed vocabulary, leaving \OovEnPercent\% of English and
\OovEsPercent\% of Spanish test tokens unknown.

\maybefigure{data_subwords}{1.0}{What segmentation costs. Subword tokenisation
lengthens sequences by roughly a third, which is real compute, in exchange for
never encountering an unknown word.}{fig:subwords}

\subsection{Batching}

Sequences are padded to a rectangle within each batch, and attention masks mark
which positions are real. Batches are formed by \emph{token count} rather than
sentence count: the length distribution is right-skewed, so a fixed sentence
count would produce batches whose padding waste swings widely depending on
whether a long sentence happened to be drawn. The sampler shuffles, takes pools
of fifty batches' worth of examples, sorts each pool by length, and emits
batches capped at 8{,}192 tokens. Batch order is then shuffled again so the
model does not systematically see short sentences first within an epoch.

% ---------------------------------------------------------------------------
\section{Model architecture}
\label{sec:architecture}

\maybefigure{transformer_architecture}{0.92}{The complete model, annotated with
the shapes this project actually uses. Read bottom to top on each side. The
encoder turns the tagged source into one vector per source position; the
decoder generates the target one token at a time, attending to what it has
already produced and, through cross-attention, to the encoder's
output.}{fig:architecture}

\subsection{Attention}

The single operation the architecture is built around is scaled dot-product
attention:
\begin{equation}
\mathrm{Attention}(Q, K, V) =
  \mathrm{softmax}\!\left(\frac{QK^{\top}}{\sqrt{\dk}}\right)V .
\label{eq:attention}
\end{equation}

Read it as a \emph{differentiable dictionary lookup}. Each query vector is
compared against every key by dot product; the scores become a probability
distribution through softmax; the output is that distribution's weighted
average of the value vectors. A hard lookup would take the single best-matching
key --- attention takes a soft blend, which is what makes it differentiable and
therefore trainable end to end.

\paragraph{Why divide by $\sqrt{\dk}$.} If the components of $q$ and $k$ are
independent with mean~0 and variance~1, then $q \cdot k = \sum_{i=1}^{\dk} q_i
k_i$ has mean~0 and variance $\dk$, so its standard deviation grows as
$\sqrt{\dk}$. With $\dk = 64$ the raw scores would routinely reach $\pm 8$.
Softmax over values that spread out saturates: it becomes nearly one-hot, and
the gradient it passes back is proportional to $p(1-p) \approx 0$. Dividing by
$\sqrt{\dk}$ restores unit variance and is the difference between a model that
trains and one that stalls.

\maybefigure{attention_mechanism}{1.0}{Scaled dot-product attention as a chain
of matrix operations, with the shapes used here.}{fig:attention}

\paragraph{Why multiple heads.} A single softmax produces one distribution per
query, so one head can attend to essentially one place at a time. Translating
``the red house'' into ``la casa roja'' requires simultaneously tracking the
noun (for gender agreement on both the article and the adjective) and the
adjective (to move it after the noun). Splitting $\dmodel$ into $h$ independent
subspaces of size $\dmodel/h$ lets $h$ such relations be attended to in
parallel at identical total cost --- the split is a reshape, not extra
computation.

\maybefigure{multi_head_split}{0.95}{Multi-head attention: $h$ parallel
subspaces at one head's cost.}{fig:heads}

\subsection{Positional encoding}

Self-attention is permutation-equivariant: it computes a weighted sum, and a
sum does not care about order. Shuffle the input words and every attention
output is shuffled identically --- the model literally cannot distinguish ``the
dog bit the man'' from ``the man bit the dog''. Position must be injected
explicitly. We use fixed sinusoidal encodings,
\begin{align}
PE_{(pos,\,2i)}   &= \sin\!\left(pos / 10000^{2i/\dmodel}\right), \\
PE_{(pos,\,2i+1)} &= \cos\!\left(pos / 10000^{2i/\dmodel}\right),
\end{align}
whose motivating property is that \emph{relative} position is a linear function
of absolute position: for any fixed offset $k$ there is a matrix $M_k$ --- a
$2\times 2$ rotation applied per frequency pair --- with $PE_{pos+k} = M_k\,
PE_{pos}$. Attention can therefore learn ``attend three tokens back'' as one
linear map that works at every position, and because the encoding is a function
rather than a lookup table it extends to lengths never seen in training. A
learned alternative is implemented for the ablation.

\maybefigure{positional_encoding}{1.0}{The sinusoidal table and the waves that
generate it. Wavelengths run from $2\pi$ to $10000\cdot 2\pi$: fast dimensions
separate neighbouring words, slow ones separate distant
regions.}{fig:positional}

\subsection{Masking}

Two logically different masks are needed, and conflating them is a classic
source of silent bugs. The \emph{padding} mask hides filler positions added to
square off the batch; without it, the encoder's representation of a short
sentence would depend on how long the other sentences in its batch happened to
be. The \emph{causal} mask stops the decoder from seeing the future: during
training the whole target is supplied at once so all positions compute in
parallel, but position $t$ must attend only to positions $\le t$. Without it
the model trivially copies the next token from its input, scores near-perfect
training loss, and is useless at inference.

Throughout the codebase \texttt{True} means ``attend here''. The alternative
convention is equally common, and mixing the two silently inverts the model.

\maybefigure{masking}{1.0}{The two masks and their combination for decoder
self-attention. A key is usable only if it is both real and not in the
future.}{fig:masking}

\subsection{Residual connections and layer normalisation}

Each sublayer is wrapped in a residual connection around a normalised
transformation. The residual gives the gradient a path back to the input that
does not pass through the sublayer: differentiating $x + f(x)$ gives $1 +
f'(x)$, so even if $f'$ vanishes the $1$ survives. It also changes what each
block must learn --- only a \emph{correction} to a running representation,
rather than the whole representation.

We use \textbf{pre-layer-normalisation}, $x \leftarrow x + \mathrm{Sublayer}
(\mathrm{LayerNorm}(x))$, rather than the original paper's post-LN. In pre-LN
the residual path runs from input to output un-normalised, so gradients reach
the early layers undamped and training is stable at higher learning rates with
a shorter warmup. Post-LN needs a carefully tuned warmup to avoid diverging in
the first few hundred steps. The practical argument is decisive for a project
trained in free Colab sessions: a run that diverges at step 300 is an hour
wasted. Pre-LN requires one extra \code{LayerNorm} at the end of each stack,
since the last sublayer's output would otherwise never be normalised.

\maybefigure{residual_norm}{0.92}{Where to put the layer normalisation, and why
it changes training stability.}{fig:prenorm}

\subsection{Hyper-parameter choices}
\label{sec:hyperparameters}

\begin{table}[htbp]
\centering
\caption{The systems compared. All four share the same data, direction tags,
optimiser, decoding code and evaluation.}
\label{tab:model}
\small
\begin{tabular}{lrrlr}
\toprule
System & Vocab & Emb. dim & Geometry & Params \\
\midrule
BPE 16k, learned embeddings & 16,000 & 512 & 4+4 layers, $d_{\mathrm{model}}=512$, 8 heads & 37.6M \\
Word 32k, random init & 32,000 & 300 & 4+4 layers, $d_{\mathrm{model}}=512$, 8 heads & 39.4M \\
Word 32k, MUSE init & 32,000 & 300 & 4+4 layers, $d_{\mathrm{model}}=512$, 8 heads & 39.4M \\
LSTM + Bahdanau attention & 16,000 & 512 & 2 layers, hidden 704, bi-encoder & 38.9M \\
\bottomrule
\end{tabular}

\end{table}

\begin{decisionbox}
\textbf{$\dmodel = 512$, $h = 8$, $\dff = 2048$.} These are the proportions of
``Transformer base'' \cite{vaswani2017}: $\dff = 4\dmodel$ and $\dk =
\dmodel/h = 64$. Staying on well-understood ratios means the published guidance
about learning rates and warmup transfers directly, so the experimental budget
can be spent on the questions this project is actually asking.

\textbf{4 encoder + 4 decoder layers}, reduced from the paper's 6+6. WMT14
English--German has roughly 4.5M sentence pairs; this corpus has
\PairsTrain{}, two orders of magnitude less. Depth is the first thing to cut
when data rather than compute is the binding constraint, and halving it also
halves step time and memory, which is what makes an epoch fit comfortably in a
free Colab session.

\textbf{Dropout 0.1}, the paper's value, applied to sublayer outputs and to the
attention weights. Dropout on the \emph{weights} rather than the outputs is
what the original specifies: it forces the model to route information through
alternative positions, so no single alignment becomes load-bearing.

\textbf{Three-way tied embeddings.} The encoder embedding, the decoder
embedding and the output projection are the same matrix, which is possible only
because of the joint vocabulary. At $\dmodel = 512$ and $V = \SubwordVocab$
this removes about 17M parameters --- roughly a third of the model. On a corpus
this size that is regularisation as much as economy: the alternative is three
separate tables each seeing a third of the gradient signal.
\end{decisionbox}

\subsection{Decoupling embedding width from model width}

The pre-trained-embedding experiments use 300-dimensional MUSE vectors against
a 512-wide residual stream. Rather than shrinking $\dmodel$ to 300 --- which
would change the architecture and confound the comparison --- the model inserts
a learned $300 \rightarrow 512$ projection on the way in and a $512 \rightarrow
300$ projection on the way out, with the output logits computed against the
transposed embedding matrix. Everything between the two projections is
byte-for-byte identical across experiments, so a difference in results is
attributable to the embeddings rather than to a change of architecture.

% ---------------------------------------------------------------------------
\section{Training pipeline}
\label{sec:training}

The optimisation loop is written by hand, as the brief requires: forward, loss,
backward, clip, step, schedule, log. No high-level training utility is used.

\subsection{Loss: label-smoothed cross entropy}

Plain cross entropy asks the model to put probability 1.0 on the reference
token and 0.0 on all others. For translation that target is false --- several
outputs are genuinely correct --- and the only way to drive the loss down is to
make the logit gap enormous, which produces overconfident, badly calibrated
output. Label smoothing \cite{szegedy2016} replaces the target with
\begin{equation}
q(k) =
\begin{cases}
1 - \varepsilon + \dfrac{\varepsilon}{V'} & k = y, \\[8pt]
\dfrac{\varepsilon}{V'} & k \neq y,
\end{cases}
\end{equation}
with $\varepsilon = 0.1$. Two implementation details matter and are both easy
to get wrong: \code{<pad>} positions must be excluded from the loss \emph{and}
from the token count used to normalise it, otherwise the reported loss depends
on how much padding happened to be in the batch; and \code{<pad>} must also be
excluded from the smoothing distribution, since spreading probability mass onto
a token the model must never emit teaches it to emit that token. Hence $V' = V
- 1$.

Label smoothing deliberately raises the reported cross entropy --- the smoothed
target has non-zero entropy, so the loss cannot reach zero even for a perfect
model. We therefore log the \emph{unsmoothed} negative log-likelihood alongside
it, and compute perplexity from that, so the number stays interpretable.

\subsection{Optimiser and schedule}

AdamW with $\beta = (0.9, 0.98)$ and $\epsilon = 10^{-9}$, the paper's
settings. The higher second-moment $\beta$ (0.98 rather than the usual 0.999)
makes the optimiser adapt faster to the changing gradient scale of the warmup
phase. Weight decay of 0.01 is applied \emph{only} to genuine weight matrices:
decaying LayerNorm gains and bias terms is a common mistake and a harmful one,
since shrinking a normalisation gain towards zero suppresses the signal the
layer exists to pass through.

The learning rate follows the inverse-square-root schedule,
\begin{equation}
lr(t) = \dmodel^{-0.5} \cdot
        \min\!\left(t^{-0.5},\; t \cdot t_{\text{warmup}}^{-1.5}\right),
\end{equation}
with 4{,}000 warmup steps. The reason warmup is needed is structural: at
initialisation the attention weights are near-uniform, so every position's
output is roughly the average of all positions and carries almost no
information about \emph{which} position mattered. The resulting gradients are
dominated by noise, and Adam's second-moment estimate is itself unreliable in
the first few dozen steps --- so a large step taken then is large in an
essentially arbitrary direction. The classic symptom is a loss that falls for
200 steps and then diverges to NaN.

This schedule is not parameterised by the \emph{total} step count, so a run can
be stopped early or extended without invalidating it --- which matters when
training in Colab sessions that may be interrupted.

\maybefigure{train_lr_bpe_scratch}{1.0}{The learning-rate schedule as realised
during training, and the analytic curve behind it.}{fig:lr}

\subsection{Stability and efficiency}

\begin{itemize}
  \item \textbf{Gradient clipping} by global norm at 1.0. Transformer gradients
    are usually well behaved but occasionally spike, and one unclipped spike
    can undo an epoch of progress.
  \item \textbf{Gradient accumulation} lets a large effective batch be
    simulated on a small GPU. The 8{,}192-token batch is near the limit of a
    free Colab T4; accumulating over two or four micro-batches reaches the
    16k--32k tokens that stabilise transformer training without more memory.
  \item \textbf{Mixed precision} on CUDA, roughly halving memory and speeding
    up matrix multiplication. It is deliberately \emph{disabled} on Apple's MPS
    backend, where at this model size it is slower than fp32 and occasionally
    produces NaNs in the attention softmax.
  \item \textbf{Early stopping} on validation loss, with both the best-loss and
    the best-BLEU checkpoints retained.
\end{itemize}

\subsection{Monitoring}

Validation BLEU is computed on a fixed 500-sentence sample after every epoch,
using greedy decoding --- the purpose is a comparable \emph{trend}, not the best
achievable score. This matters because validation loss and BLEU do not peak at
the same epoch: loss typically starts rising while BLEU is still improving, so
selecting a checkpoint on loss alone leaves quality on the table.

Every metric is appended to a JSON Lines file as it is produced, so an
interrupted Colab session still leaves a complete, plottable history behind.

\maybefigure{train_curves_bpe_scratch}{1.0}{Training and validation curves for
the primary model. Loss and BLEU are shown on separate axes rather than as a
dual-axis chart: two $y$-scales on one frame invite comparison of slopes that
share no units, and the apparent crossing point would be an artefact of
scaling.}{fig:curves}

\begin{pending}
Discussion of the observed curves --- where the model begins to overfit, the
gap between the loss-optimal and BLEU-optimal epochs, and the effect of the
early-stopping patience --- is written once the full training runs are
complete. The figure and every number quoted around it regenerate from
\code{artifacts/results/} automatically.
\end{pending}

% ---------------------------------------------------------------------------
\section{Evaluation and baseline comparison}
\label{sec:evaluation}

\subsection{Metrics}

BLEU \cite{papineni2002} scores a hypothesis by how much of its $n$-gram
content appears in the reference:
\begin{equation}
\mathrm{BLEU} = BP \cdot \exp\!\left(\sum_{n=1}^{4} \tfrac{1}{4} \log p_n\right),
\qquad
BP = \begin{cases} 1 & c > r \\ e^{1 - r/c} & c \le r \end{cases}
\end{equation}
where $p_n$ is the \emph{clipped} modified precision --- each reference
$n$-gram can be matched only as many times as it occurs there, which is what
stops ``the the the the'' from scoring perfect unigram precision. Counts are
pooled over the whole corpus before the ratio is taken. Since BLEU has no
recall term, the brevity penalty stands in for one.

We report \textbf{sacreBLEU} \cite{post2018} figures as the headline numbers,
because BLEU is notoriously sensitive to tokenisation and the same system can
differ by several points depending on how text was split before counting. The
reproducibility signature is recorded with every evaluation. We also implement
BLEU from the definition and assert agreement with sacreBLEU on every
evaluation run; a unit test pins the agreement to within 0.1 BLEU, and on the
verification corpora the two implementations agree exactly.

\textbf{chrF2} \cite{popovic2015} is reported alongside. It is a character
$n$-gram F-score, and it matters specifically for Spanish: a translation can get
a word's stem right and its inflection wrong, which BLEU counts as a total miss
and chrF gives partial credit for. The gap between the two metrics is itself
informative.

\subsection{Baselines}

Reporting an absolute BLEU number in isolation tells a reader very little, so
the transformer is compared against:

\begin{enumerate}
  \item \textbf{A recurrent sequence-to-sequence model} --- bidirectional LSTM
    encoder, LSTM decoder, additive (Bahdanau) attention \cite{bahdanau2015}
    --- matched on vocabulary, data, optimiser, label smoothing, decoding code
    and parameter count. The hidden size was chosen so the baseline carries
    slightly \emph{more} parameters than the transformer, making the comparison
    conservative. Its schedule differs deliberately: the inverse-sqrt/warmup
    recipe is tuned for transformers, and applying it to an LSTM would produce
    a straw-man baseline.
  \item \textbf{Two embedding-initialisation conditions} that differ only in
    the contents of the embedding matrix at step 0 (Section~\ref{sec:muse}).
\end{enumerate}

\subsection{Decoding}

Greedy decoding takes the highest-probability token at each step: fast, and the
right default for interactive use, but myopic --- a token that looks best
locally may leave no good continuation, and greedy cannot revise it. Beam
search keeps the $k$ best partial hypotheses and extends all of them.

Two details of the beam implementation matter. A hypothesis' score is a sum of
log-probabilities, all negative, so longer hypotheses score worse purely for
being longer; unnormalised beam search therefore produces systematically
truncated translations. Dividing by $\text{length}^{0.6}$ \cite{wu2016}
corrects this. And a hypothesis that has emitted \code{</s>} must be set aside
rather than extended, or it keeps accumulating mass and crowds the beam.

\maybefigure{beam_search}{0.95}{Beam search over the target vocabulary, with
the length-normalisation issue that makes it work.}{fig:beam}

\subsection{Results}

\begin{table}[htbp]
\centering
\caption{Held-out test-set results (\PairsTest{} sentence pairs per direction,
beam search with width 4 and length penalty 0.6). BLEU and chrF2 are
sacreBLEU figures.}
\label{tab:results}
\small
\begin{tabular}{lrrrrr}
\toprule
& \multicolumn{2}{c}{EN $\rightarrow$ ES} & \multicolumn{2}{c}{ES $\rightarrow$ EN} & \\
\cmidrule(lr){2-3}\cmidrule(lr){4-5}
System & BLEU & chrF2 & BLEU & chrF2 & Mean BLEU \\
\midrule
BPE 16k, learned embeddings & 49.99 & 68.50 & 56.83 & 71.43 & 53.41 \\
Word 32k, random init & 45.42 & 65.57 & 52.62 & 68.22 & 49.02 \\
Word 32k, MUSE init & 45.88 & 65.60 & 53.05 & 68.57 & 49.46 \\
LSTM + Bahdanau attention & 46.15 & 65.71 & 49.73 & 66.37 & 47.94 \\
\bottomrule
\end{tabular}

\end{table}

\maybefigure{eval_bleu}{1.0}{Translation quality by system and
direction.}{fig:bleu}

\maybefigure{eval_length}{1.0}{Sentence-level BLEU stratified by source length.
This figure tests the central architectural claim: a recurrent model must carry
information across a number of sequential steps proportional to sentence
length, whereas in a transformer every pair of positions is one attention hop
apart, so the recurrent model should degrade faster as sentences
lengthen.}{fig:length}

\begin{pending}
Quantitative discussion of Table~\ref{tab:results} --- the size of the
transformer's margin over the recurrent baseline, whether the two directions
differ, and how the length-stratified curves compare --- is written once the
full training runs are complete.
\end{pending}

\subsection{The pre-trained embedding experiment}
\label{sec:muse}

The question is whether initialising the shared embedding table from
cross-lingual word vectors helps on a corpus of this size. The experiment is
constructed so that the answer means something.

\begin{keypoint}
\textbf{Why MUSE and not monolingual fastText.} The model has \emph{one} shared
embedding table serving both languages on both the encoder and the decoder
side. Concatenating two independently-trained monolingual tables would place
``cat'' and ``gato'' in unrelated coordinate systems; the initialisation would
be worse than random, because the model would first have to \emph{unlearn} the
accidental geometry. MUSE \cite{conneau2018} maps both languages into a shared
space through a learned orthogonal transform, so translation-equivalent words
start close together. That is a genuine prior for a shared table.
\end{keypoint}

\begin{keypoint}
\textbf{Why a third run is necessary.} Comparing the MUSE model against the
primary BPE model would confound two variables at once --- subword versus word
tokenisation, and pre-trained versus random initialisation --- and would prove
nothing. Run~B (word-level, random) exists purely as the control: it is
byte-for-byte identical to Run~C except for the contents of the embedding
matrix at step~0.
\end{keypoint}

The initialisation covers \MuseCoveragePercent\% of the word vocabulary
(\MuseCovered{} types found, \MuseRandom{} drawn at random). Uncovered entries
are sampled at the \emph{same scale} as the pre-trained cloud rather than from
the default $\mathcal{N}(0,1)$; otherwise unknown words would begin with
vectors roughly thirty times longer than known ones and would dominate the
attention logits. Pre-trained vectors are held frozen for two epochs so the
randomly-initialised encoder and decoder can adapt to the given geometry before
large early gradients scramble it, then unfrozen for fine-tuning.

That the space really is cross-lingual at initialisation can be checked
directly --- the nearest neighbours of \emph{house} include \emph{casa}, of
\emph{water} include \emph{agua}, and of \emph{book} include \emph{libro},
before a single training step.

\begin{pending}
The outcome of the comparison --- whether the MUSE initialisation helps, and
whether any advantage survives to convergence or merely accelerates early
epochs --- is written once the full training runs are complete.
\end{pending}

% ---------------------------------------------------------------------------
\section{Error analysis}
\label{sec:errors}

Reading eleven thousand test translations by hand is not practical, so the
analysis is automated: every test sentence is scored with sentence-level BLEU,
tagged with the failure modes it exhibits, and ranked. The twenty worst cases
per direction are reproduced in full in Appendix~\ref{app:examples} so that
each classification can be checked by eye --- the detectors are a \emph{reading
aid}, not a verdict.

\subsection{Failure modes detected}

\begin{description}[leftmargin=1.6cm,style=nextline]
  \item[repetition] An $n$-gram repeated more often than in the reference. The
    classic degenerate-decoding loop.
  \item[truncation / over-generation] Output much shorter or longer than the
    reference. Truncation usually means an early \code{</s>}; over-generation
    often accompanies repetition.
  \item[unknown\_token] The output contains \code{<unk>}. Possible only for the
    word-level models, and quantifying this gap is much of the point of the
    subword comparison.
  \item[copied\_source] The output overlaps the \emph{source} more than the
    reference: the model gave up and passed the input through. Frequent with
    rare proper nouns.
  \item[number\_mismatch] Digits in the reference are missing or altered.
    Separated out because it is a \emph{semantic} error that BLEU
    under-penalises, and the kind of mistake that makes a translation system
    unusable in practice regardless of its score.
  \item[no\_content\_overlap] Not a single content word shared with the
    reference --- either a complete failure or a legitimate paraphrase BLEU
    cannot see. Both are worth inspecting, and the appendix distinguishes them
    by hand.
\end{description}

\begin{table}[htbp]
\centering
\caption{Failure-mode frequency across the test set for \PrimaryRun. Categories
are not mutually exclusive: one translation can exhibit several.}
\label{tab:errors}
\small
\begin{tabular}{lrr}
\toprule
Failure mode & EN $\rightarrow$ ES & ES $\rightarrow$ EN \\
\midrule
\texttt{other} & 94.7\% & 96.3\% \\
\texttt{no\_content\_overlap} & 3.8\% & 2.3\% \\
\texttt{repetition} & 0.6\% & 0.6\% \\
\texttt{truncation} & 0.4\% & 0.4\% \\
\texttt{over\_generation} & 0.4\% & 0.3\% \\
\texttt{number\_mismatch} & 0.2\% & 0.4\% \\
\texttt{copied\_source} & 0.2\% & 0.1\% \\
\bottomrule
\end{tabular}

\end{table}

\maybefigure{eval_errors_bpe_scratch}{1.0}{Failure-mode census for the primary
model.}{fig:errors}

\begin{pending}
The discussion of which failure modes dominate, what causes them, and which are
addressable --- together with the manual reading of the twenty worst examples
per direction --- is written once the full training runs are complete.
\end{pending}

% ---------------------------------------------------------------------------
\section{Inference application}
\label{sec:app}

The trained model is served through a Gradio application
(\code{app/app.py}) that loads a saved checkpoint once at start-up, accepts any
sentence the user types, translates it in either direction without retraining,
and displays the output immediately.

Beyond the minimum, the interface exposes the parts of the system this report
discusses, so that a viewer can interrogate the model rather than only use it:

\begin{itemize}
  \item a \textbf{direction switch} --- the same weights serve both ways, and
    the app changes only the direction tag prepended to the source;
  \item \textbf{decoding controls} --- greedy against beam search, beam width
    and length penalty, so the effects described in
    Section~\ref{sec:evaluation} can be reproduced live;
  \item an \textbf{attention view} showing the cross-attention alignment for
    the sentence just translated;
  \item a \textbf{tokenisation view}, which is what makes an unexpected
    translation of a rare word explicable.
\end{itemize}

The checkpoint carries its own model configuration and the path of the
tokeniser it was trained with, so the application needs only a file path and
can serve any of the four systems without code changes.

\maybefigure{eval_decoding}{0.85}{What the search contributes: BLEU under
greedy decoding against beam search. A large gap means the model's single-best
token is often wrong even when its distribution is
right.}{fig:decoding}

% ---------------------------------------------------------------------------
\section{Limitations and future work}

\begin{itemize}
  \item \textbf{Domain.} Tatoeba is short, simple, everyday sentences. The
    system should not be expected to handle technical, legal or literary text,
    and the length-stratified results in Figure~\ref{fig:length} show quality
    falling on the longest inputs, which are also the rarest in training.
  \item \textbf{Single reference.} BLEU is computed against one reference per
    sentence, so legitimate paraphrases are penalised. Tatoeba's many-to-many
    links could supply multiple references, at the cost of a more complex
    evaluation set --- a natural extension.
  \item \textbf{No key/value cache.} The decoder re-runs over the whole prefix
    at each generation step, $O(n^2)$ work where $O(n)$ would do. This was a
    deliberate trade: the cache would roughly double the size of the decoder
    code and obscure the architecture this report explains, and at a median of
    eight subword tokens the wall-clock cost is negligible. It would matter for
    longer inputs.
  \item \textbf{Capacity shared between directions.} A single model of a given
    size must hold both directions; two specialised models would likely score
    higher each, at twice the parameters and half the supervision per
    parameter.
  \item \textbf{Back-translation} is the obvious next step for a corpus this
    size: monolingual Spanish and English text is abundant, and synthetic
    parallel data from a reverse model is the standard way to exploit it.
\end{itemize}

% ---------------------------------------------------------------------------
\section{Conclusion}

We built a complete neural machine translation system --- corpus construction,
preprocessing, a transformer implemented from primitives, a hand-written
training loop, evaluation against a matched baseline, automated error analysis
and an interactive application --- for English and Spanish in both directions
with a single set of weights.

Two methodological points are worth carrying away independently of the scores.
First, the way a corpus is split can matter more than the model: Tatoeba's
graph structure makes a naive random split leak, and the resulting BLEU
measures memorisation. Second, an experiment comparing pre-trained embeddings
against random initialisation is only interpretable if the tokenisation is held
fixed, which is why this project trains three models rather than two.

% ---------------------------------------------------------------------------
\begin{thebibliography}{9}
\small

\bibitem{vaswani2017}
A.~Vaswani, N.~Shazeer, N.~Parmar, J.~Uszkoreit, L.~Jones, A.~N. Gomez,
  {\L}.~Kaiser, and I.~Polosukhin.
\newblock Attention is all you need.
\newblock In \emph{Advances in Neural Information Processing Systems 30}, 2017.

\bibitem{johnson2017}
M.~Johnson et al.
\newblock Google's multilingual neural machine translation system: enabling
  zero-shot translation.
\newblock \emph{Transactions of the ACL}, 5:339--351, 2017.

\bibitem{bahdanau2015}
D.~Bahdanau, K.~Cho, and Y.~Bengio.
\newblock Neural machine translation by jointly learning to align and
  translate.
\newblock In \emph{ICLR}, 2015.

\bibitem{sennrich2016}
R.~Sennrich, B.~Haddow, and A.~Birch.
\newblock Neural machine translation of rare words with subword units.
\newblock In \emph{ACL}, 2016.

\bibitem{kudo2018}
T.~Kudo and J.~Richardson.
\newblock SentencePiece: a simple and language independent subword tokenizer
  and detokenizer for neural text processing.
\newblock In \emph{EMNLP: System Demonstrations}, 2018.

\bibitem{conneau2018}
A.~Conneau, G.~Lample, M.~Ranzato, L.~Denoyer, and H.~J{\'e}gou.
\newblock Word translation without parallel data.
\newblock In \emph{ICLR}, 2018.

\bibitem{papineni2002}
K.~Papineni, S.~Roukos, T.~Ward, and W.-J. Zhu.
\newblock BLEU: a method for automatic evaluation of machine translation.
\newblock In \emph{ACL}, 2002.

\bibitem{post2018}
M.~Post.
\newblock A call for clarity in reporting BLEU scores.
\newblock In \emph{WMT}, 2018.

\bibitem{popovic2015}
M.~Popovi{\'c}.
\newblock chrF: character $n$-gram F-score for automatic MT evaluation.
\newblock In \emph{WMT}, 2015.

\bibitem{szegedy2016}
C.~Szegedy, V.~Vanhoucke, S.~Ioffe, J.~Shlens, and Z.~Wojna.
\newblock Rethinking the Inception architecture for computer vision.
\newblock In \emph{CVPR}, 2016.

\bibitem{wu2016}
Y.~Wu et al.
\newblock Google's neural machine translation system: bridging the gap between
  human and machine translation.
\newblock \emph{arXiv:1609.08144}, 2016.

\bibitem{xiong2020}
R.~Xiong et al.
\newblock On layer normalization in the transformer architecture.
\newblock In \emph{ICML}, 2020.

\bibitem{tatoeba}
The Tatoeba Project.
\newblock \url{https://tatoeba.org/en/downloads}. Accessed 2026.

\end{thebibliography}

% ---------------------------------------------------------------------------
\appendix
\clearpage
\section{Reproducing this work}
\label{app:repro}

\begin{lstlisting}[language=bash]
git clone https://github.com/jolusano/en-es-transformer-nmt.git
cd en-es-transformer-nmt
python -m venv .venv && source .venv/bin/activate
pip install -r requirements.txt

python -m nmt.data.build                                  # corpus + tokenisers
python -m nmt.data.embeddings                             # MUSE init matrix
python -m nmt.training.train --config configs/bpe_scratch.yaml
python -m nmt.evaluation.evaluate \
    --checkpoint artifacts/checkpoints/bpe_scratch/best_bleu.pt

python -m nmt.viz.make_figures                            # all figures
python reports/build_report_data.py                       # all tables
python app/app.py                                         # the application
\end{lstlisting}

The full instructions, including the Colab route for training, are in the
repository \code{README.md}.

\clearpage
\section{The twenty worst translations per direction}
\label{app:examples}

Ranked by sentence-level BLEU, from the primary model (\PrimaryRun). Sentence
BLEU is meaningful only for ranking sentences against one another; it is not
comparable to the corpus figures in Table~\ref{tab:results}.

\small
\subsection*{EN $\rightarrow$ ES}
\begin{longtable}{@{}rp{0.7cm}p{11cm}@{}}
\toprule
\# & sBLEU & Source / reference / model output / failure modes \\
\midrule
\endhead
1 & 0.0 & \small Tell me about it! \\
& & \small\itshape ref: Ni me lo digas. \\
& & \small\ttfamily out: ¡Cuéntame sobre ello! \\
& & \scriptsize no\_content\_overlap \\\addlinespace
2 & 0.0 & \small Sit down! \\
& & \small\itshape ref: ¡Sentaos! \\
& & \small\ttfamily out: Siéntese. \\
& & \scriptsize no\_content\_overlap \\\addlinespace
3 & 0.0 & \small Curse you! \\
& & \small\itshape ref: Te maldigo. \\
& & \small\ttfamily out: ¡Durmentaos! \\
& & \scriptsize no\_content\_overlap \\\addlinespace
4 & 0.0 & \small That's really great! \\
& & \small\itshape ref: Es realmente magnífico. \\
& & \small\ttfamily out: ¡Qué bien es eso! \\
& & \scriptsize no\_content\_overlap \\\addlinespace
5 & 0.0 & \small Sit down! \\
& & \small\itshape ref: ¡Sentate! \\
& & \small\ttfamily out: Siéntese. \\
& & \scriptsize no\_content\_overlap \\\addlinespace
6 & 0.0 & \small Tell me about it. \\
& & \small\itshape ref: ¡Dímelo a mí! \\
& & \small\ttfamily out: Háblame de ello. \\
& & \scriptsize no\_content\_overlap \\\addlinespace
7 & 0.0 & \small Sit down! \\
& & \small\itshape ref: ¡Siéntate! \\
& & \small\ttfamily out: Siéntese. \\
& & \scriptsize no\_content\_overlap \\\addlinespace
8 & 3.0 & \small Wealthy older men often marry younger trophy wives. \\
& & \small\itshape ref: Hombres ricos suelen casarse con jóvenes mujeres trofeo. \\
& & \small\ttfamily out: A menudo, los hombres mayores, a menudo nos casan con las trofeas más jóvenes. \\
& & \scriptsize over\_generation \\\addlinespace
9 & 3.4 & \small A strong argument for the religion of Christ is this - that offences against Charity are about the only ones which men on their death-beds can be made, not to understand, but to feel, as crime. \\
& & \small\itshape ref: Un sólido argumento en favor del cristianismo es el siguiente: las ofensas contra la caridad es probablemente lo único que, en sus lechos de muerte, los hombres llegan a sentir y no a comprender como un crimen. \\
& & \small\ttfamily out: Un argumento fuerte para la religión de Cristo es esta - que las oficiales de Charualidad están a punto de los únicos que los hombres de su muerte pueden ser obligados, no para entender, sino que no se sienten el crimen. \\
& & \scriptsize other \\\addlinespace
10 & 3.5 & \small The squadron encountered an ambush and scrambled for coverage. \\
& & \small\itshape ref: El escuadrón se topó con una emboscada y se apresuró a ponerse a cubierto. \\
& & \small\ttfamily out: El caspa encontró un embosco de ambush y scorrocaba por cobertura. \\
& & \scriptsize no\_content\_overlap \\\addlinespace
11 & 3.6 & \small His meaning is quite plain. \\
& & \small\itshape ref: Lo que él trata de decir es bien sencillo de comprender. \\
& & \small\ttfamily out: Su significado es bastante simple. \\
& & \scriptsize truncation, no\_content\_overlap \\\addlinespace
12 & 3.7 & \small Peasants often have a secondary activity to augment their income. \\
& & \small\itshape ref: A menudo los campesinos se ocupan en una actividad secundaria para complementar sus ingresos. \\
& & \small\ttfamily out: Los Peas a menudo tienen una segunda actividad mayor a un aumento de sueldo. \\
& & \scriptsize other \\\addlinespace
13 & 3.8 & \small Mr. Smith makes it a rule to take a walk every morning. \\
& & \small\itshape ref: El Señor Smith tiene el hábito de salir a caminar todas las mañanas. \\
& & \small\ttfamily out: El señor Smith hace como norma dar un paseo cada mañana. \\
& & \scriptsize other \\\addlinespace
14 & 4.0 & \small You can go out on condition that you come home by seven. \\
& & \small\itshape ref: Si vuelves antes de las 7, puedes salir. \\
& & \small\ttfamily out: Puedes salir con la condición de que vengas a casa a las siete. \\
& & \scriptsize number\_mismatch \\\addlinespace
15 & 4.1 & \small What marvelous work he's doing, donating books to needy children! \\
& & \small\itshape ref: ¡Qué maravilloso trabajo que está haciendo en donar libros a los niños necesitados! \\
& & \small\ttfamily out: ¡Qué trabajo tan maravilloso es hacer, no hacer libros que necesitan tener hijos! \\
& & \scriptsize other \\\addlinespace
16 & 4.1 & \small Galician, Portuguese, and Spanish are often confused with one another. \\
& & \small\itshape ref: A menudo se confunden el gallego, el portugués y el español. \\
& & \small\ttfamily out: Galician, portugués, y español a menudo son confundidos con uno al otro. \\
& & \scriptsize other \\\addlinespace
17 & 4.2 & \small The electric bill keeps going up. \\
& & \small\itshape ref: Cada vez sube más el precio de la electricidad. \\
& & \small\ttfamily out: La cuenta eléctrica sigue subiendo. \\
& & \scriptsize no\_content\_overlap \\\addlinespace
18 & 4.4 & \small We make a teeny bit of progress, then we go back to square one. \\
& & \small\itshape ref: Apenas logramos un avance, todo vuelve a cero. \\
& & \small\ttfamily out: Hacemos un poco de progreso, entonces volvimos a cuadrarnos a una. \\
& & \scriptsize no\_content\_overlap \\\addlinespace
19 & 4.5 & \small It's coming out! My bladder's about to explode! \\
& & \small\itshape ref: ¡Me estoy meando, la vejiga me va a explotar! \\
& & \small\ttfamily out: ¡Ya está fuera! ¡Mi vejiga está a punto de estallar! \\
& & \scriptsize other \\\addlinespace
20 & 4.5 & \small My little sister has been suffering with a toothache since last night. \\
& & \small\itshape ref: A mi hermana menor le duele un diente desde esta noche. \\
& & \small\ttfamily out: Mi hermana pequeña lleva con dolor de muelas desde anoche. \\
& & \scriptsize other \\\addlinespace
\bottomrule
\end{longtable}

\subsection*{ES $\rightarrow$ EN}
\begin{longtable}{@{}rp{0.7cm}p{11cm}@{}}
\toprule
\# & sBLEU & Source / reference / model output / failure modes \\
\midrule
\endhead
1 & 0.0 & \small Aquí están diez dólares. \\
& & \small\itshape ref: Here's \$10.00. \\
& & \small\ttfamily out: Ten dollars are here. \\
& & \scriptsize number\_mismatch, no\_content\_overlap \\\addlinespace
2 & 0.0 & \small Te maldigo. \\
& & \small\itshape ref: Curse you! \\
& & \small\ttfamily out: I swear. \\
& & \scriptsize no\_content\_overlap \\\addlinespace
3 & 3.1 & \small Considera a las mujeres placeres de usar y tirar y no búsquedas con sentido. \\
& & \small\itshape ref: He regards women as disposable pleasures rather than as meaningful pursuits. \\
& & \small\ttfamily out: Consider the women of pleasure to use and throw away and don't make sense. \\
& & \scriptsize other \\\addlinespace
4 & 3.2 & \small Considera que las mujeres son placeres de usar y tirar más que búsquedas con sentido. \\
& & \small\itshape ref: He regards women as disposable pleasures rather than as meaningful pursuits. \\
& & \small\ttfamily out: It seems that women are delights to use and throw more than dinner with meaning. \\
& & \scriptsize other \\\addlinespace
5 & 3.7 & \small Considera a las mujeres placeres desechables y no propósitos con sentido. \\
& & \small\itshape ref: He regards women as disposable pleasures rather than as meaningful pursuits. \\
& & \small\ttfamily out: Consider the women with disposal women and do not propose with meaning. \\
& & \scriptsize other \\\addlinespace
6 & 4.0 & \small ¿Rusos nuevos, dices? ¡Maravilloso! ¡Es justo lo que necesitábamos! \\
& & \small\itshape ref: New Russians, you say? Wonderful! That's just what we need! \\
& & \small\ttfamily out: You're new, you're saying? What we needed! \\
& & \scriptsize other \\\addlinespace
7 & 4.1 & \small Un sólido argumento en favor del cristianismo es el siguiente: las ofensas contra la caridad es probablemente lo único que, en sus lechos de muerte, los hombres llegan a sentir y no a comprender como un crimen. \\
& & \small\itshape ref: A strong argument for the religion of Christ is this - that offences against Charity are about the only ones which men on their death-beds can be made, not to understand, but to feel, as crime. \\
& & \small\ttfamily out: A solid argument in favor of the Christianity is the next: the offensive against unity is probably the only thing in his deaths, men came to feel and not understanding as a crime. \\
& & \scriptsize other \\\addlinespace
8 & 4.1 & \small Estados Unidos se imagina que es la nación más libre del mundo. \\
& & \small\itshape ref: The United States fancies itself the world's freest nation. \\
& & \small\ttfamily out: America is imagining that it is the free nation in the world. \\
& & \scriptsize other \\\addlinespace
9 & 4.4 & \small Estados Unidos se imagina que es la nación más libre del mundo. \\
& & \small\itshape ref: America fancies itself the world's freest nation. \\
& & \small\ttfamily out: America is imagining that it is the free nation in the world. \\
& & \scriptsize other \\\addlinespace
10 & 4.5 & \small No apresure el paso; llegaremos a la hora. \\
& & \small\itshape ref: Don't walk so fast. We'll get there on time. \\
& & \small\ttfamily out: Don't rush the step; we'll arrive at the hour. \\
& & \scriptsize other \\\addlinespace
11 & 4.5 & \small Considera que las mujeres son placeres desechables más que objetivos significativos. \\
& & \small\itshape ref: He regards women as disposable pleasures rather than as meaningful pursuits. \\
& & \small\ttfamily out: I consider women to be desirable more than mean goals. \\
& & \scriptsize other \\\addlinespace
12 & 4.7 & \small Los baños de los aviones no tienen ventanillas por motivos de seguridad. \\
& & \small\itshape ref: For security reasons, the bathrooms on airplanes don't have windows. \\
& & \small\ttfamily out: Airplanes have no windows for safety reasons. \\
& & \scriptsize other \\\addlinespace
13 & 4.8 & \small Cada vez sube más el precio de la electricidad. \\
& & \small\itshape ref: The electric bill keeps going up. \\
& & \small\ttfamily out: Every time it rises the price of electricity. \\
& & \scriptsize no\_content\_overlap \\\addlinespace
14 & 4.8 & \small No hay falsedad en lo que él dice. \\
& & \small\itshape ref: What he's saying is true. \\
& & \small\ttfamily out: There are no falsehood in what he says. \\
& & \scriptsize other \\\addlinespace
15 & 4.8 & \small Como suele pasar, él no me trajo nada. \\
& & \small\itshape ref: As usual he brought me nothing. \\
& & \small\ttfamily out: As is often the case, he didn't bring me anything. \\
& & \scriptsize over\_generation \\\addlinespace
16 & 4.9 & \small Ella puso en riesgo su vida para salvar a un niño de ahogarse. \\
& & \small\itshape ref: She saved the drowning child at the risk of her own life. \\
& & \small\ttfamily out: She risked her life to save a child from drowning. \\
& & \scriptsize other \\\addlinespace
17 & 4.9 & \small Tenga cuidado de no resbalar en las baldosas mojadas. \\
& & \small\itshape ref: Mind you don't slip on the wet tiles. \\
& & \small\ttfamily out: Be careful not to slippery. \\
& & \scriptsize no\_content\_overlap \\\addlinespace
18 & 4.9 & \small Por favor, ajuste su asiento como le acomode. \\
& & \small\itshape ref: Please adjust the seat to fit you. \\
& & \small\ttfamily out: Please put your seat on as a hold of him. \\
& & \scriptsize other \\\addlinespace
19 & 4.9 & \small Lo sentimos mucho por no poder ayudarlos. \\
& & \small\itshape ref: We are sorry we can't help you. \\
& & \small\ttfamily out: We're very sorry for not being able to help them. \\
& & \scriptsize other \\\addlinespace
20 & 4.9 & \small ¿Cuál fue el momento en el que te diste cuenta de que habías crecido? \\
& & \small\itshape ref: When did you realize you've grown up? \\
& & \small\ttfamily out: What was the time you noticed that you had grown? \\
& & \scriptsize other \\\addlinespace
\bottomrule
\end{longtable}


\end{document}
